% =====================================================================
%  THÈME 4 — Logarithmes, pH & incertitude — EXERCICES — Niveau DIFFICILE
% =====================================================================
\documentclass[11pt,a4paper]{article}
\usepackage[utf8]{inputenc}
\usepackage[T1]{fontenc}
\usepackage{lmodern}
\usepackage[french]{babel}
\usepackage{amsmath,amssymb,mathtools}
\usepackage{icomma}
\usepackage{siunitx}
\sisetup{output-decimal-marker={,},per-mode=symbol}
\usepackage[version=4]{mhchem}
\usepackage{xcolor}
\usepackage[most]{tcolorbox}
\usepackage{enumitem}
\usepackage{fancyhdr}
\usepackage[a4paper,margin=2.2cm,headheight=26pt,headsep=8pt]{geometry}

\definecolor{primary}{RGB}{150,98,162}
\definecolor{primarydark}{RGB}{92,52,110}
\newcommand{\cs}{chiffre significatif}
\newcommand{\CS}{chiffres significatifs}

\pagestyle{fancy}\fancyhf{}
\fancyhead[L]{\small\textcolor{primarydark}{\textbf{Fiche 1} — Chiffres significatifs}}
\fancyhead[R]{\small\textcolor{primarydark}{\textsc{Exercices · Difficile · Thème 4}}}
\fancyfoot[C]{\small\thepage}
\renewcommand{\headrulewidth}{0.8pt}
\renewcommand{\headrule}{\hbox to\headwidth{\color{primary}\leaders\hrule height \headrulewidth\hfill}}

\newtcolorbox[auto counter]{exo}[1][]{enhanced,breakable,boxrule=0.8pt,arc=2pt,
  colback=primary!4,colframe=primary,left=3mm,right=3mm,top=2mm,bottom=2mm,
  fonttitle=\bfseries,coltitle=white,
  title={Exercice~\thetcbcounter\ifx\relax#1\relax\relax\else\ ---\ \textmd{\itshape #1}\fi}}

\setlength{\parindent}{0pt}\setlength{\parskip}{3pt}

\begin{document}

\begin{center}
{\Large\bfseries\color{primarydark} Exercices — Niveau Difficile}\\[1pt]
{\color{primary} Thème 4 : pH \& incertitude — synthèse}
\end{center}
\vspace{1mm}
\textit{On combine la règle du logarithme et l'analyse de la précision. Traite les sous-questions dans
l'ordre.}
\vspace{2mm}

\begin{exo}[pH dans les deux sens, et chiffres réellement significatifs]
Une solution vérifie $[\ce{H3O+}]=\num{3,0e-5}\,\si{\mol\per\litre}$. On donne $\log(\num{3,0})=\num{0,48}$
et $10^{\num{0,48}}=\num{3,0}$.
\begin{enumerate}[label=\alph*),leftmargin=1.6em,itemsep=1pt]
  \item Calcule le pH avec le bon nombre de décimales.
  \item Parmi les chiffres du pH obtenu, lesquels sont \emph{significatifs} ? Que représente la partie
        entière ?
  \item Vérifie ta réponse en repartant du pH pour retrouver $[\ce{H3O+}]$ (avec le bon nombre de CS).
\end{enumerate}
\end{exo}

\begin{exo}[Deux balances : beaucoup de CS ne veut pas dire « juste »]
On pèse un étalon de masse \emph{vraie} $\qty{2,000}{\gram}$. La balance A affiche, de façon
reproductible, $\qty{2,0500}{\gram}$ ; la balance B affiche, de façon reproductible, $\qty{2,00}{\gram}$.
\begin{enumerate}[label=\alph*),leftmargin=1.6em,itemsep=1pt]
  \item Donne l'incertitude absolue puis l'incertitude relative (en \%) de chaque affichage (d'après
        le dernier \cs).
  \item Quelle balance est la plus \textbf{fidèle} (précise à l'écriture) ?
  \item Quelle balance est la plus \textbf{juste} (proche de la vraie valeur) ? Justifie par l'écart
        à $\qty{2,000}{\gram}$.
  \item Un grand nombre de \CS\ peut-il compenser un défaut de justesse ? Conclus.
\end{enumerate}
\end{exo}

\begin{exo}[Un titrage : de l'incertitude au pH]
Un titrage donne $[\ce{H3O+}]=(\num{2,0}\pm\num{0,1})\times10^{-3}\,\si{\mol\per\litre}$. On donne
$\log(\num{2,0})=\num{0,30}$.
\begin{enumerate}[label=\alph*),leftmargin=1.6em,itemsep=1pt]
  \item Calcule l'incertitude \emph{relative} sur $[\ce{H3O+}]$ (en \%). Combien de \CS\ ce niveau de
        précision justifie-t-il pour la concentration ?
  \item Calcule le pH avec le bon nombre de décimales.
  \item Explique, en une phrase, le lien entre le nombre de CS de $[\ce{H3O+}]$ et le nombre de
        décimales du pH.
\end{enumerate}
\end{exo}

\end{document}
